% ===================================================================== PREAMBULE COMMUN — Carnet d'entraînement Fiche 9
\documentclass[11pt,a4paper]{article}
\usepackage[utf8]{inputenc}
\usepackage[T1]{fontenc}
\usepackage{lmodern}
\usepackage[french]{babel}
\usepackage{amsmath,amssymb,mathtools}
\usepackage{icomma}
\usepackage{siunitx}
\sisetup{output-decimal-marker={,},per-mode=symbol}
\DeclareSIUnit{\molar}{M}
\DeclareSIUnit{\Molar}{mol\,L^{-1}}
\usepackage[version=4]{mhchem}
\usepackage{xcolor}
\usepackage{colortbl}
\usepackage{tikz}
\usepackage{pgfplots}
\usepackage[most]{tcolorbox}
\usepackage{enumitem}
\usepackage{array}
\usepackage{booktabs}
\usepackage{multicol}
\usepackage{fancyhdr}
\usepackage[hidelinks]{hyperref}
\usetikzlibrary{arrows.meta,positioning,calc,decorations.pathmorphing,
  decorations.markings,angles,quotes,patterns,backgrounds,
  shapes.geometric,shapes.misc,fadings,intersections,fit}
\pgfplotsset{compat=1.18}
\usepackage[a4paper,margin=2.2cm,headheight=26pt,headsep=10pt]{geometry}

% ---------- Palette ----------
\definecolor{primary}{RGB}{150,98,162}
\definecolor{primarydark}{RGB}{92,52,110}
\definecolor{defcol}{RGB}{0,114,178}
\definecolor{thmcol}{RGB}{0,140,120}
\definecolor{formulacol}{RGB}{198,93,9}
\definecolor{remarkcol}{RGB}{120,94,200}
\definecolor{attncol}{RGB}{200,40,55}
\definecolor{excol}{RGB}{0,135,90}
\definecolor{methcol}{RGB}{90,90,100}
\definecolor{cationcol}{RGB}{0,90,190}
\definecolor{anioncol}{RGB}{200,60,55}
\definecolor{cristalcol}{RGB}{110,105,120}
\definecolor{eaucol}{RGB}{120,180,220}
\definecolor{solcol}{RGB}{0,135,90}
\definecolor{kpscol}{RGB}{198,93,9}
\definecolor{qcol}{RGB}{120,94,200}
\definecolor{precipcol}{RGB}{110,70,40}
\definecolor{exercol}{RGB}{150,98,162}
\definecolor{corrcol}{RGB}{0,120,100}

% ---------- Macros ----------
\newcommand{\Kps}{K_{\mathrm{ps}}}
\newcommand{\Ce}[1]{\left[\,\ce{#1}\,\right]}

% ---------- Style des boîtes ----------
\tcbset{boxbase/.style={enhanced,breakable,boxrule=0.9pt,arc=2.5pt,
   left=3mm,right=3mm,top=2.3mm,bottom=2.3mm,
   fonttitle=\bfseries,coltitle=white,toptitle=0.6mm,bottomtitle=0.6mm}}

\newtcolorbox[auto counter]{definition}[1][]{%
  boxbase,colback=defcol!5,colframe=defcol,
  title={\textsc{Définition}~\thetcbcounter\ \ifx\relax#1\relax\relax\else\textmd{--- \itshape #1}\fi}}
\newtcolorbox[auto counter]{theoreme}[1][]{%
  boxbase,colback=thmcol!6,colframe=thmcol,
  title={\textsc{Loi / Propriété}~\thetcbcounter\ \ifx\relax#1\relax\relax\else\textmd{--- \itshape #1}\fi}}
\newtcolorbox[auto counter]{formule}[1][]{%
  boxbase,colback=formulacol!6,colframe=formulacol,
  title={\textsc{Formule}~\thetcbcounter\ \ifx\relax#1\relax\relax\else\textmd{--- \itshape #1}\fi}}
\newtcolorbox{remarque}[1][]{%
  boxbase,colback=remarkcol!6,colframe=remarkcol,
  title={\textsc{Remarque}\ifx\relax#1\relax\relax\else\textmd{ : \itshape #1}\fi}}
\newtcolorbox{attention}[1][]{%
  boxbase,colback=attncol!6,colframe=attncol,
  title={\textsc{Attention}\ifx\relax#1\relax\relax\else\textmd{ : \itshape #1}\fi}}
\newtcolorbox{exemple}[1][]{%
  boxbase,colback=excol!5,colframe=excol,
  title={\textsc{Exemple}\ifx\relax#1\relax\relax\else\textmd{ : \itshape #1}\fi}}
\newtcolorbox{methode}[1][]{%
  boxbase,colback=methcol!7,colframe=methcol,sharp corners,
  title={\textsc{Méthode}\ifx\relax#1\relax\relax\else\textmd{ : \itshape #1}\fi}}
\newtcolorbox{astuce}[1][]{%
  boxbase,colback=primary!7,colframe=primary,
  title={\textsc{Astuce}\ifx\relax#1\relax\relax\else\textmd{ : \itshape #1}\fi}}

\newtcolorbox[auto counter]{exercice}[1][]{%
  boxbase,colback=exercol!4,colframe=exercol,
  title={\textsc{Exercice}~\thetcbcounter\ \ifx\relax#1\relax\relax\else\textmd{--- \itshape #1}\fi}}
\newtcolorbox[auto counter]{correction}[1][]{%
  boxbase,colback=corrcol!4,colframe=corrcol,
  title={\textsc{Corrigé}~\thetcbcounter\ \ifx\relax#1\relax\relax\else\textmd{--- \itshape #1}\fi}}

% ---------- Environnement « où : » ----------
\newenvironment{ou}{%
  \par\smallskip\noindent\textbf{où :}\nopagebreak
  \begin{itemize}[leftmargin=1.8em,topsep=2pt,itemsep=1.5pt,parsep=0pt]}%
  {\end{itemize}}

% ---------- Styles TikZ ----------
\tikzset{
  cat/.style={circle,fill=cationcol,draw=cationcol!50!black,inner sep=0pt,
              minimum size=5mm,font=\bfseries\scriptsize,text=white},
  an/.style={circle,fill=anioncol,draw=anioncol!50!black,inner sep=0pt,
             minimum size=5mm,font=\bfseries\scriptsize,text=white},
  crx/.style={circle,fill=cristalcol,draw=black!50,inner sep=0pt,
              minimum size=5mm,font=\bfseries\scriptsize,text=white},
  ptn/.style={circle,fill=black,inner sep=1.1pt}}

% ---------- En-têtes ----------
\pagestyle{fancy}
\fancyhf{}
\fancyhead[L]{\small\textcolor{primarydark}{\textbf{Fiche 9} — Solubilité et produit de solubilité}}
\fancyhead[R]{\small\textcolor{primarydark}{\textsc{Carnet d'entraînement}}}
\fancyfoot[C]{\small\thepage}
\renewcommand{\headrulewidth}{0.8pt}
\renewcommand{\headrule}{\hbox to\headwidth{\color{primary}\leaders\hrule height \headrulewidth\hfill}}
\usepackage{titlesec}
\titleformat{\section}{\Large\bfseries\color{primarydark}}{\thesection}{0.6em}{}
\titleformat{\subsection}{\large\bfseries\color{primary}}{\thesubsection}{0.6em}{}

% ---------- Bandeau de titre ----------
% #1 = thème/étiquette   #2 = titre principal   #3 = sous-titre
\newcommand{\bandeau}[3]{%
\begin{center}
\begin{tikzpicture}
\node[rectangle,rounded corners=7pt,fill=primary,text=white,
      minimum width=15.6cm,minimum height=1.95cm,align=center,inner sep=8pt]
  {{\large\bfseries #1}\\[3pt]{\Large\bfseries #2}\\[3pt]{\small #3}};
\end{tikzpicture}
\end{center}
\vspace{2mm}}
\begin{document}
\bandeau{Thème 4 — Effet d'ion commun}%
{Exercices — Niveau Difficile}%
{Récupération quantitative, courbe $s'(C)$, et le piège de l'ion commun cationique}

\noindent\textit{Consigne : sous-questions progressives. Attention : la formule $s'=\Kps/C^{\,n}$
n'est valable que si l'ion commun est l'\emph{anion} $\mathrm{X}^-$ — sinon, il faut la redémontrer.}
\medskip

\begin{exercice}[récupérer l'argent d'une solution]
On veut « chasser » l'argent d'une solution en le faisant précipiter sous forme de \ce{AgCl}
($\Kps=\num{1{,}8e-10}$). On travaille sur \qty{1{,}0}{\litre} de solution. Donnée : $M(\ce{Ag})=\qty{108}{\gram\per\mole}$.
\begin{enumerate}[label=\textbf{\alph*)},leftmargin=2.2em,topsep=2pt,itemsep=2pt]
  \item Rappelle la solubilité $s_0$ de \ce{AgCl} dans l'eau pure, et la quantité de matière d'argent
        dissous (en mol) dans \qty{1{,}0}{\litre} saturé.
  \item On ajoute des chlorures jusqu'à $[\ce{Cl-}]=\qty{0{,}10}{\molar}$. Calcule la nouvelle
        solubilité $s'=[\ce{Ag+}]$.
  \item Combien de moles d'argent restent \emph{dissoutes} dans \qty{1{,}0}{\litre} ? Quelle masse
        d'argent cela représente-t-il ?
  \item Quelle \textbf{fraction} de l'argent initialement dissous a été précipitée ? Commente
        l'efficacité du procédé.
\end{enumerate}
\end{exercice}

\begin{exercice}[courbe de solubilité $s'$ en fonction de l'ion commun]
Pour \ce{CaF2} ($\Kps=\num{3{,}9e-11}$), on étudie la solubilité $s'$ en fonction de la concentration
$C=[\ce{F-}]$ imposée. La figure ci-dessous trace $s'(C)$ en échelle log–log.

\begin{center}
\begin{tikzpicture}
\begin{axis}[
   width=10.5cm,height=5.4cm,
   xmode=log,ymode=log,
   xlabel={$C=[\ce{F-}]$ imposée (\si{\molar})},
   ylabel={$s'$ (\si{\molar})},
   xmin=1e-3,xmax=1,ymin=1e-11,ymax=1e-4,
   axis lines=left,grid=both,grid style={black!12},
   every axis x label/.style={at={(xticklabel cs:1)},anchor=west},
   every axis y label/.style={at={(yticklabel cs:1)},anchor=south,rotate=-90}]
  \addplot[thick,anioncol,mark=*,mark options={fill=anioncol}] coordinates {
    (1e-3,3.9e-5) (1e-2,3.9e-7) (1e-1,3.9e-9) (1,3.9e-11)};
\end{axis}
\end{tikzpicture}
\end{center}

\begin{enumerate}[label=\textbf{\alph*)},leftmargin=2.2em,topsep=2pt,itemsep=2pt]
  \item Établis l'expression de $s'$ en fonction de $\Kps$ et $C$ (type $1{:}2$, ion commun \ce{F-}).
  \item Calcule $s'$ pour $C=\qty{0{,}050}{\molar}$ puis pour $C=\qty{0{,}20}{\molar}$.
  \item Sur la courbe log–log, la droite descend : quand $C$ est multipliée par 10, par combien $s'$
        est-elle divisée ? Relie ce facteur à l'exposant de $C$ dans $s'=\Kps/C^{2}$.
\end{enumerate}
\end{exercice}

\begin{exercice}[le piège : ion commun \emph{cationique} sur un sel $1{:}2$]
On dissout du bromure de plomb \ce{PbBr2} ($\Kps=\num{9{,}2e-6}$) dans une solution de nitrate de plomb
\ce{Pb(NO3)2} telle que $[\ce{Pb^2+}]=\qty{0{,}10}{\molar}$. Ici l'ion commun est le \textbf{cation}
\ce{Pb^2+}.
\begin{enumerate}[label=\textbf{\alph*)},leftmargin=2.2em,topsep=2pt,itemsep=2pt]
  \item Un camarade applique directement $s'=\Kps/C^{2}$. Explique pourquoi c'est \textbf{faux} ici
        (quel ion vient de la dissolution, quel ion est imposé ?).
  \item Exprime $[\ce{Br-}]$ en fonction de la solubilité $s'$, puis écris $\Kps$ en fonction de $s'$ et
        de $C=[\ce{Pb^2+}]$.
  \item Isole $s'$. Tu dois obtenir $s'=\dfrac{1}{2}\sqrt{\dfrac{\Kps}{C}}$.
  \item Calcule $s'$ et compare à la solubilité $s_0=\qty{1{,}3e-2}{\molar}$ de \ce{PbBr2} dans l'eau pure.
\end{enumerate}
\end{exercice}

\end{document}
